\documentclass[11pt,a4paper,fontset=fandol]{ctexart}

\usepackage[a4paper,top=2.25cm,bottom=2.45cm,left=2.25cm,right=2.25cm,headheight=18pt,headsep=0.55cm,footskip=24pt]{geometry}
\usepackage{amsmath,amssymb,amsthm}
\usepackage{fancyhdr}
\usepackage{lastpage}
\usepackage{enumitem}
\usepackage{titlesec}
\usepackage{xcolor}
\usepackage{hyperref}
\usepackage{bookmark}
\usepackage[most]{tcolorbox}
\usepackage{setspace}
\usepackage{array}

\hypersetup{
  colorlinks=true,
  linkcolor=blue!50!black,
  urlcolor=blue!50!black
}

\setstretch{1.13}
\setlength{\parindent}{2em}
\setlength{\parskip}{0.25em}
\allowdisplaybreaks
\setlist[enumerate]{itemsep=0.45em, topsep=0.35em, leftmargin=2.4em}
\setlist[itemize]{itemsep=0.25em, topsep=0.25em, leftmargin=1.9em}

\definecolor{mainblue}{RGB}{36,83,140}
\definecolor{lightblue}{RGB}{241,246,252}
\definecolor{lightgray}{RGB}{246,246,246}
\definecolor{rulegray}{RGB}{110,118,128}

\titleformat{\section}
  {\Large\bfseries\color{mainblue}}
  {\thesection.}
  {0.45em}
  {}
  [\vspace{0.2em}\color{rulegray}\titlerule]
\titlespacing*{\section}{0pt}{1.1em}{0.7em}

\pagestyle{fancy}
\fancyhf{}
\fancyhead[L]{\small 容斥原理周测 B 卷}
\fancyhead[C]{\small 组合专题：综合变式与错位结合}
\fancyhead[R]{\small 刘文博}
\fancyfoot[C]{\small 第 \thepage 页 / 共 \pageref{LastPage} 页}
\renewcommand{\headrulewidth}{0.55pt}
\renewcommand{\footrulewidth}{0.35pt}

\newtcolorbox{infobox}[1]{
  breakable,
  colback=lightblue,
  colframe=mainblue,
  boxrule=0.7pt,
  arc=1mm,
  title=\textbf{#1},
  fonttitle=\bfseries
}

\newenvironment{solution}{\par\noindent\textbf{参考解答：}}{\hfill$\square$\par}
\newcommand{\answerspace}{\vspace{2.3cm}}
\newcommand{\biganswerspace}{\vspace{3.2cm}}
\newcommand{\Dn}{D_n}

\begin{document}

\thispagestyle{empty}
\begin{center}
{\fontsize{22}{28}\selectfont\bfseries 容斥原理周测 B 卷\par}
\vspace{0.4cm}
{\large 适合小学高年级至初中低年级数学竞赛训练\par}
\end{center}

\begin{infobox}{考试说明}
\begin{itemize}
  \item 测试时间：75 分钟
  \item 满分：100 分
  \item B 卷更强调模型迁移，请特别注意：什么时候按“至少一个条件”正面算，什么时候按“坏事件并集”来处理。
\end{itemize}
\end{infobox}

\section{试题}

\begin{enumerate}
  \item （12 分）在 $1$ 到 $100$ 中，是 2 的倍数或 3 的倍数，但不是 5 的倍数的整数有多少个？
  \answerspace

  \item （12 分）7 个不同对象排成一行，要求前 4 个对象都不在原位，而后 3 个对象没有限制。共有多少种排法？
  \biganswerspace

  \item （12 分）8 个不同对象排成一行，要求前 6 个对象都不在原位，后 2 个对象没有限制。共有多少种排法？
  \biganswerspace

  \item （12 分）7 封不同的信放入 7 个对应信封中，至少有 2 封装对，共有多少种装法？
  \answerspace

  \item （12 分）用容斥原理计算 $D_6$。
  \answerspace

  \item （12 分）8 个不同数字排成一行，至少有 3 个数字在原位，共有多少种排法？
  \biganswerspace

  \item （14 分）9 个不同对象排成一行。要求前 4 个对象中至少有 1 个在原位。共有多少种排法？
  \biganswerspace

  \item （14 分）证明：若共有 $n$ 个对象，其中前 $m$ 个对象都不能回到原位，而其余对象没有限制，则满足条件的排列数为
  \[
  \sum_{k=0}^{m}(-1)^k\binom{m}{k}(n-k)!.
  \]
  \biganswerspace
\end{enumerate}

\newpage
\section{详细解答}

\begin{enumerate}
  \item \begin{solution}
  先求 1 到 100 中“是 2 的倍数或 3 的倍数”的整数个数。

  设
  \[
  A=\{\text{2 的倍数}\},\qquad B=\{\text{3 的倍数}\}.
  \]

  则
  \[
  |A\cup B|=50+33-16=67.
  \]

  再求这些数里同时还是 5 的倍数的个数，也就是“10 的倍数或 15 的倍数”的个数。

  设
  \[
  C=\{\text{10 的倍数}\},\qquad D=\{\text{15 的倍数}\}.
  \]

  则
  \[
  |C|=10,
  \qquad
  |D|=6,
  \qquad
  |C\cap D|=\left\lfloor\frac{100}{30}\right\rfloor=3.
  \]

  所以这部分共有
  \[
  10+6-3=13
  \]
  个。

  因此所求为
  \[
  67-13=54.
  \]
  \end{solution}

  \item \begin{solution}
  设前 4 个对象中，第 $i$ 个回到原位的事件为 $A_i$。

  总排法有
  \[
  7!=5040
  \]
  种。

  单个坏事件共有
  \[
  \binom41 6!=4\times 720=2880
  \]
  种。

  两个坏事件同时发生共有
  \[
  \binom42 5!=6\times 120=720
  \]
  种。

  三个坏事件同时发生共有
  \[
  \binom43 4!=4\times 24=96
  \]
  种。

  四个坏事件同时发生共有
  \[
  \binom44 3!=6
  \]
  种。

  所以所求为
  \[
  5040-2880+720-96+6=2790.
  \]
  \end{solution}

  \item \begin{solution}
  设前 6 个对象回到原位的坏事件为 $A_1,A_2,\ldots,A_6$。

  总排法有
  \[
  8!=40320
  \]
  种。

  单个坏事件共有
  \[
  \binom61 7!=6\times 5040=30240
  \]
  种。

  两个坏事件同时发生共有
  \[
  \binom62 6!=15\times 720=10800
  \]
  种。

  三个坏事件同时发生共有
  \[
  \binom63 5!=20\times 120=2400
  \]
  种。

  四个坏事件同时发生共有
  \[
  \binom64 4!=15\times 24=360
  \]
  种。

  五个坏事件同时发生共有
  \[
  \binom65 3!=6\times 6=36
  \]
  种。

  六个坏事件同时发生共有
  \[
  \binom66 2!=2
  \]
  种。

  所以所求为
  \[
  40320-30240+10800-2400+360-36+2=18806.
  \]
  \end{solution}

  \item \begin{solution}
  “至少有 2 封装对”按“恰好有几封装对”分类。

  \textbf{恰好有 2 封装对：}
  \[
  \binom72 D_5=21\times 44=924.
  \]

  \textbf{恰好有 3 封装对：}
  \[
  \binom73 D_4=35\times 9=315.
  \]

  \textbf{恰好有 4 封装对：}
  \[
  \binom74 D_3=35\times 2=70.
  \]

  \textbf{恰好有 5 封装对：}
  \[
  \binom75 D_2=21.
  \]

  \textbf{恰好有 6 封装对：}
  不可能。

  \textbf{恰好有 7 封装对：}
  有 1 种。

  所以总数为
  \[
  924+315+70+21+1=1331.
  \]
  \end{solution}

  \item \begin{solution}
  设 $A_i$ 表示“第 $i$ 个对象在原位”。则
  \[
  D_6=6!-\binom61 5!+\binom62 4!-\binom63 3!+\binom64 2!-\binom65 1!+\binom66 0!.
  \]

  代入得
  \[
  D_6=720-720+360-120+30-6+1=265.
  \]

  因此
  \[
  D_6=265.
  \]
  \end{solution}

  \item \begin{solution}
  “至少有 3 个在原位”按“恰好有几个在原位”分类。

  \textbf{恰好有 3 个在原位：}
  \[
  \binom83 D_5=56\times 44=2464.
  \]

  \textbf{恰好有 4 个在原位：}
  \[
  \binom84 D_4=70\times 9=630.
  \]

  \textbf{恰好有 5 个在原位：}
  \[
  \binom85 D_3=56\times 2=112.
  \]

  \textbf{恰好有 6 个在原位：}
  \[
  \binom86 D_2=28.
  \]

  \textbf{恰好有 7 个在原位：}
  不可能。

  \textbf{恰好有 8 个在原位：}
  有 1 种。

  所以总数为
  \[
  2464+630+112+28+1=3235.
  \]
  \end{solution}

  \item \begin{solution}
  “前 4 个对象中至少有 1 个在原位”适合用补集法。

  总排法有
  \[
  9!=362880
  \]
  种。

  先求前 4 个对象一个都不在原位的排法数。设这 4 个对象回原位的事件为 $A_1,A_2,A_3,A_4$，则这一部分为
  \[
  9!-\binom41 8!+\binom42 7!-\binom43 6!+\binom44 5!.
  \]

  逐项计算：
  \[
  4\times 8!=161280,
  \qquad
  6\times 7!=30240,
  \]
  \[
  4\times 6!=2880,
  \qquad
  5!=120.
  \]

  所以前 4 个对象一个都不在原位的排法数为
  \[
  362880-161280+30240-2880+120=229080.
  \]

  因此所求为
  \[
  362880-229080=133800.
  \]
  \end{solution}

  \item \begin{solution}
  设 $A_i$ 表示“前 $m$ 个受限制对象中的第 $i$ 个回到原位”。

  我们要计算的是“这些坏事件一个都不发生”的排列数，也就是
  \[
  \overline{A_1\cup A_2\cup \cdots\cup A_m}
  \]
  的大小。

  总排法有
  \[
  n!
  \]
  种。

  如果指定恰有 $k$ 个受限制对象回到原位，那么这 $k$ 个位置就固定住了，剩余还有
  \[
  n-k
  \]
  个对象可以任意排列，因此对应排法数为
  \[
  (n-k)!.
  \]

  这样的 $k$ 个受限制对象有
  \[
  \binom{m}{k}
  \]
  种选法。

  于是由容斥原理，所求为
  \[
  n!-\binom{m}{1}(n-1)!+\binom{m}{2}(n-2)!-\cdots+(-1)^m\binom{m}{m}(n-m)!.
  \]

  这正是
  \[
  \sum_{k=0}^{m}(-1)^k\binom{m}{k}(n-k)!.
  \]

  证毕。
  \end{solution}
\end{enumerate}

\end{document}
